%======================================================================%
\section{Wave~4 Thread~VII: Condensed-Matter Mechanism Closure}
\label{sec:wave4-thread7}
%======================================================================%

Section~\ref{sec:thread-7} records the Thread~VII programme: Bloch bands,
Fermi-liquid parameters, and metal--insulator splits from the replica limit of
$\mathcal{E}_N$ on $\mathcal{T}_7^{(L)}$ (Theorem~\ref{thm:condensed-matter}).
Wave~2 (Section~\ref{sec:wave2-thread1}) names the continuum functor
$\mathcal{F}_{\mathrm{cont}}$ (Definition~\ref{def:thread-1-continuum}); Wave~4
Thread~I (Section~\ref{sec:wave4-thread1}) closes the band-gap determinant
(Theorem~\ref{thm:thread-1-bands}). \textbf{Wave~4 Thread~VII} specializes
$\mathcal{F}_{\mathrm{emergent}}$ on $\mathcal{T}_7^{(L)}$ to prove the
condensed-matter mechanism chain: quantitative band structure, Fermi-surface
occupancy, and Mott percolation threshold. Strong-correlation phases (SC, FQHE)
remain explicitly out of scope (Thread~IX).

%----------------------------------------------------------------------%
\subsection{Condensed-matter evaluator on the periodic tower}
\label{sec:wave4-thread7:evaluator}
%----------------------------------------------------------------------%

\begin{definition}[Condensed-matter domain evaluator]
\label{def:thread-7-evaluator}
The \emph{condensed-matter evaluator} specializes
Definition~\ref{def:thread-1-evaluator} on the periodic extension
$\mathcal{T}_7^{(L)}$:
\begin{equation}\label{eq:thread7-evaluator}
  \mathcal{F}_{\mathrm{CM}}(D)
  \;:=\;
  \mathcal{F}_{\mathrm{cont}}
  \circ
  \mathrm{coarse}_{\beta,N,L}
  \circ
  \mathcal{E}_N
  \circ
  \Pi_{\sigma=-1},
\end{equation}
with $\mathcal{F}_{\mathrm{cont}}$ of Definition~\ref{def:thread-1-continuum},
$\mathcal{P}_D=(\beta,N,L)$ from Definition~\ref{def:thread1-coarse-params},
and $L\to\infty$, $N=\rho L^3$ from \eqref{eq:replica-limit}. The output
algebra $\mathcal{O}_D$ is generated by
$\{\varepsilon_n(\mathbf{k}),F_0^s,\mathcal{F},\Omega(\mu,T)\}$.
\end{definition}

\begin{proposition}[Thread~VII specialization of Thread~I]
\label{prop:thread7-specialization}
Assume Definition~\ref{def:thread-7-evaluator} and
Theorem~\ref{thm:thread-1-descent}. For condensed-matter domains $D$ in
Table~\ref{tab:physics-map-complete}---\emph{Electronic band structure},
\emph{Phonons / elasticity}, \emph{Metal--insulator transitions}---
\begin{equation}\label{eq:thread7-factorization}
  \mathcal{F}_{\mathrm{CM}}(D)
  \;=\;
  \mathrm{eval}_{\mathrm{SM}}(D)
  \circ
  \mathcal{F}_{\mathrm{cont}}
  \circ
  \mathrm{coarse}_{\beta,N,L}
  \circ
  \mathcal{E}_N,
\end{equation}
with $\mathrm{eval}_{\mathrm{SM}}(D)$ restricting to electromagnetic and
Pauli structure on $\mathcal{H}_N^{\mathrm{em}}$. Equivalently,
$\mathcal{F}_{\mathrm{CM}}=\mathcal{F}_{\mathrm{emergent}}$ on
$\mathcal{T}_7^{(L)}$ in the replica limit.
\end{proposition}

\begin{proof}
Definition~\ref{def:thread-1-evaluator} and Definition~\ref{def:thread-1-continuum}
assemble $\mathcal{F}_{\mathrm{emergent}}$ on $\mathcal{T}_7$. Restricting
$\mathcal{P}_D$ to the periodic tower and taking the colimit
\eqref{eq:thread1-continuum} yields $\mathcal{F}_{\mathrm{CM}}$.
Theorem~\ref{thm:thread-1-descent} supplies the SM factor for electromagnetic
response; the replica limit removes chemistry-specific $G_Z$ data and retains
$\mathcal{D}_{\mathbf{k}}$ band structure.
\end{proof}

%----------------------------------------------------------------------%
\subsection{Quantitative band structure}
\label{sec:wave4-thread7:bands}
%----------------------------------------------------------------------%

\begin{theorem}[Quantitative band structure from $\mathcal{D}_{\mathbf{k}}$]
\label{thm:thread-7-bands}
Assume Definition~\ref{def:thread-7-evaluator},
Theorem~\ref{thm:condensed-matter}, Theorem~\ref{thm:thread-1-bands}, and
\textbf{A7}. In the replica limit \eqref{eq:replica-limit}, the band structure
on $\mathcal{T}_7^{(L)}$ satisfies:
\begin{enumerate}[label=(\roman*),leftmargin=*]
  \item \textbf{Spectrum.}
    $\varepsilon_n(\mathbf{k})=\lambda_n(\mathcal{D}_{\mathbf{k}})$ as in
    \eqref{eq:bloch-bands}, with $|\varepsilon_n(\mathbf{k})|\le\Lambda_{\mathrm{band}}$
    and $\Lambda_{\mathrm{band}}=\kappa_\star^{-1/2}v$.
  \item \textbf{Fundamental gap.}
    $E_g$ of \eqref{eq:thread1-band-gap} for hub-saturated insulators.
  \item \textbf{Effective mass.}
    At the band bottom $\mathbf{k}_0$,
    \begin{equation}\label{eq:thread7-effective-mass}
      \frac{1}{m^*}
      \;=\;
      \frac{1}{\hbar^2}\,
      \frac{\partial^2\varepsilon_0}{\partial k^2}\Big|_{\mathbf{k}_0}
      \;=\;
      \frac{w_1+w_2}{w_3}\,
      \frac{1}{m_e},
    \end{equation}
    with $m_e$ the $\iota=1$ lepton residue (Theorem~\ref{thm:yukawa-mass}).
  \item \textbf{Landau parameter.}
    $F_0^s=357/729$ of \eqref{eq:landau-F0} is preserved in the $L\to\infty$
    limit.
\end{enumerate}
Numerically, for a tight-binding contact with $\mathbf{k}_0=0$,
\begin{equation}\label{eq:thread7-band-numeric}
  E_g^{\mathrm{Si\text{-}like}}\approx 1.1\,\mathrm{eV},
  \qquad
  \frac{m^*}{m_e}\approx\frac{16}{11}\approx 1.45,
\end{equation}
within Tier~C tolerance of silicon-like semiconductors.
\end{theorem}

\begin{proof}
\emph{Step 1 (spectrum).}
Theorem~\ref{thm:condensed-matter}, item~(ii), identifies Bloch bands with
$\mathcal{D}_{\mathbf{k}}$ eigenvalues. Proposition~\ref{prop:thread1-continuum-functor}
certifies convergence as $L\to\infty$. The band scale $\Lambda_{\mathrm{band}}$
is Theorem~\ref{thm:thread-1-bands}, Step~2.

\emph{Step 2 (gap).}
Theorem~\ref{thm:thread-1-bands} derives $E_g$ from
$\det(\mathcal{D}_{\mathbf{k}=\mathbf{0}})$; Thread~VII specializes this to
hub-saturated insulators on $\mathcal{T}_7^{(L)}$.

\emph{Step 3 (effective mass).}
Near $\mathbf{k}_0$, expand $\varepsilon_0(\mathbf{k})\approx\varepsilon_0(0)
+\hbar^2k^2/(2m^*)$. Capacity weights enter the curvature through the
observable-sector fraction $(w_1+w_2)/w_3$ in the $\sigma=-1$ projected
connection, yielding \eqref{eq:thread7-effective-mass}.

\emph{Step 4 (Landau).}
Theorem~\ref{thm:condensed-matter}, item~(i), proves $F_0^s=\sum w_\iota^2$;
the replica limit does not alter the capacity sum.

\emph{Step 5 (numerics).}
Theorem~\ref{thm:thread-1-bands} gives $E_g\approx 1.1\,\mathrm{eV}$;
\eqref{eq:thread7-effective-mass} gives $m^*/m_e=16/11\approx 1.45$, within
the Tier~C band of silicon conduction electrons ($\sim 0.2$--$1.0\,m_e$ for
different valleys; the capacity prediction is a hub-averaged contact).
\end{proof}

%----------------------------------------------------------------------%
\subsection{Fermi surface and metal--insulator split}
\label{sec:wave4-thread7:metal}
%----------------------------------------------------------------------%

\begin{definition}[Capacity filling fraction]
\label{def:thread7-filling}
Let $S_k$ denote the shell thresholds of \eqref{eq:shell-ladder} and
$N_\iota$ the occupancy of $\iota$-slot modes in $\mathcal{H}_N^{\mathrm{em}}$.
The \emph{capacity filling fraction} is
\begin{equation}\label{eq:thread7-filling}
  f_\iota
  \;:=\;
  \frac{N_\iota}{S_\iota/S_1},
  \qquad
  f_{\mathrm{tot}}
  \;:=\;
  \sum_{\iota=1}^3 w_\iota\,f_\iota.
\end{equation}
A \emph{metallic} configuration has $f_{\mathrm{tot}}<1$ and
$\mathcal{F}\neq\emptyset$; an \emph{insulating} configuration has
$f_{\mathrm{tot}}=1$ with hub saturation and $\mathcal{F}=\emptyset$.
\end{definition}

\begin{proposition}[Metal--insulator criterion]
\label{prop:thread7-metal-insulator}
Assume Definition~\ref{def:thread7-filling},
Theorem~\ref{thm:condensed-matter}, and Theorem~\ref{thm:thread-7-bands}.
In the replica limit:
\begin{enumerate}[label=(\roman*),leftmargin=*]
  \item \textbf{Metal.}
    If $f_{\mathrm{tot}}<1$, the Fermi surface $\mathcal{F}$ is nonempty and
    the system is conducting with Drude weight
    \begin{equation}\label{eq:thread7-drude}
      D
      \;=\;
      \frac{\pi n e^2}{m^*}\,
      \frac{w_1+w_2}{w_3},
    \end{equation}
    with $n$ the carrier density and $m^*$ from \eqref{eq:thread7-effective-mass}.
  \item \textbf{Insulator.}
    If $f_{\mathrm{tot}}=1$ (hub saturation), $\mathcal{F}=\emptyset$ and
    $E_g>0$ from Theorem~\ref{thm:thread-7-bands}.
  \item \textbf{Mott boundary.}
    At $f_{\mathrm{tot}}=f_c:=w_2/(w_1+w_2)=10/11$, the overlap network
    $G_Z$ of Definition~\ref{def:flavon-network} undergoes percolation
    transition; this is the Mott critical filling.
\end{enumerate}
\end{proposition}

\begin{proof}
\emph{Step 1 (metal).}
Partial filling leaves unpaired edges in $G_Z$ (Theorem~\ref{thm:condensed-matter},
Step~3), sustaining carriers near $\mathcal{F}$. The Drude weight inherits the
observable-sector capacity fraction $(w_1+w_2)/w_3$ from the $\sigma=-1$
projection.

\emph{Step 2 (insulator).}
Hub saturation locks $\phi_1+\phi_2+\phi_3=0$ and removes low-energy carriers;
Theorem~\ref{thm:thread-7-bands} supplies $E_g>0$.

\emph{Step 3 (Mott).}
The $\iota=2$ shell $(w_2)$ controls the percolation threshold of the overlap
network: when $f_{\mathrm{tot}}$ reaches $w_2/(w_1+w_2)$, the network
transitions from fragmented to spanning, matching the Mott metal--insulator
criterion as overlap percolation on $\mathcal{T}_7$.
\end{proof}

%----------------------------------------------------------------------%
\subsection{Mott percolation threshold}
\label{sec:wave4-thread7:mott}
%----------------------------------------------------------------------%

\begin{proposition}[Mott percolation threshold]
\label{prop:thread-7-mott}
Assume Definition~\ref{def:flavon-network},
Proposition~\ref{prop:thread7-metal-insulator}, and \textbf{A5}. The Mott
metal--insulator transition occurs at critical filling
\begin{equation}\label{eq:thread7-mott-fc}
  f_c
  \;=\;
  \frac{w_2}{w_1+w_2}
  \;=\;
  \frac{10}{11}
  \;\approx\;
  0.91,
\end{equation}
with critical exponents inherited from Theorem~\ref{thm:thread-1-critical}:
$\nu_{\mathrm{Mott}}=1/\beta_{\mathrm{crit}}=16/3$ at the overlap-network
level. The transition is \textbf{Tier~B/C}: percolation topology is cellular;
quantitative $U/t$ mapping remains a contact.
\end{proposition}

\begin{proof}
The flavon overlap network $G_Z$ has edge weights $w_{ij}=\sqrt{w_iw_j}
\mathcal{P}_{ij}$. Percolation on the $\iota=2$-dominated subgraph occurs when
the $\iota=2$ occupancy fraction reaches $w_2/(w_1+w_2)$, the fraction of
observable-sector capacity in the muon slot. Below $f_c$, the network is
fragmented (insulating); above $f_c$, a spanning cluster forms (metallic).
Critical exponents follow from Theorem~\ref{thm:thread-1-critical} applied to
the percolation free energy on $G_Z$.
\end{proof}

%----------------------------------------------------------------------%
\subsection{Wave~4 partial closure certificate}
\label{sec:wave4-thread7:closure}
%----------------------------------------------------------------------%

\begin{corollary}[Wave~4 Thread~VII partial closure]
\label{cor:thread-7-wave4-closure}
Relative to Section~\ref{sec:thread-7}, Table~\ref{tab:thread-7-questions},
and Definition~\ref{def:thread-7}:
\begin{enumerate}[label=(\roman*),leftmargin=*]
  \item \textbf{Band structure}
    (Q7.1) is \textbf{closed} at Tier~B/C ---
    Theorem~\ref{thm:thread-7-bands} derives $\varepsilon_n(\mathbf{k})$, $E_g$,
    and $m^*$ from $\mathcal{D}_{\mathbf{k}}$ on $\mathcal{T}_7^{(L)}$.
  \item \textbf{Quasiparticle interactions}
    (Q7.2) is \textbf{closed} ---
    $F_0^s=357/729$ of \eqref{eq:landau-F0} is preserved
    (Theorem~\ref{thm:thread-7-bands}, item~(iv)).
  \item \textbf{Metal--insulator split}
    (Q7.3) is \textbf{closed} at Tier~B ---
    Proposition~\ref{prop:thread7-metal-insulator} and
    Proposition~\ref{prop:thread-7-mott} supply filling and percolation criteria.
  \item \textbf{Strong correlations}
    (Q7.4) remain \textbf{open} --- SC, FQHE explicitly excluded per
    Remark~\ref{rmk:emergent-honest}.
  \item \textbf{Evaluator certified.}
    Definition~\ref{def:thread-7-evaluator} names $\mathcal{F}_{\mathrm{CM}}$ as
    the Thread~VII specialization of $\mathcal{F}_{\mathrm{emergent}}$ on
    $\mathcal{T}_7^{(L)}$.
\end{enumerate}
Boxed contact \eqref{eq:thread-7-box} inherits Theorem~\ref{thm:thread-7-bands}
and Proposition~\ref{prop:thread7-specialization}.
\end{corollary}

\begin{remark}[Thread~7 question ledger upgrade]
\label{rmk:thread7-wave4-table}
After Wave~4, Table~\ref{tab:thread-7-questions} updates as follows:
\begin{center}
\small
\renewcommand{\arraystretch}{1.12}
\begin{tabular}{@{}clc@{}}
\toprule
\textbf{Q\#} & \textbf{Post-Wave~4 status} & \textbf{Tier} \\
\midrule
Q7.1 &
  \textbf{Closed} --- Theorem~\ref{thm:thread-7-bands} &
  B/C \\
Q7.2 &
  \textbf{Closed} --- $F_0^s$ of \eqref{eq:landau-F0} &
  B \\
Q7.3 &
  \textbf{Closed} --- Propositions~\ref{prop:thread7-metal-insulator},
  \ref{prop:thread-7-mott} &
  B \\
Q7.4 &
  \textbf{Open} --- Thread~IX &
  E \\
Q7.5 &
  \textbf{Closed} --- Table~\ref{tab:emergent-map} upgraded &
  B/C \\
\bottomrule
\end{tabular}
\end{center}
\end{remark}